\chapter{Extremes of stationary sequences}
\label{chap:dependence}

This chapter extends the limit theory of maxima from iid samples to stationary sequences, following
Leadbetter, Lindgren and Rootz\'en, \emph{Extremes and Related Properties of Random Sequences and
Processes}, Springer (1983) [LLR], Chapter~3, and [EKM, \S8.1]. Long-range independence of
maxima is the condition $D(u_n)$, and local clustering is controlled by $D'(u_n)$; the extremal
index $\theta$ measures how much extremes cluster.

\section{Stationary sequences and Leadbetter's conditions}

\begin{definition}[Stationary sequence]
  \label{def:stationary}
  \lean{ProbabilityTheory.IsStationary, ProbabilityTheory.IsStationary.map_shift, ProbabilityTheory.IsStationary.identDistrib,
    ProbabilityTheory.IsStationary.factor, ProbabilityTheory.isStationary_of_iIndepFun}
  \leanok
  A sequence $(X_t)_{t \in \bN}$ is stationary if the law of the shifted path $(X_{t+1})_t$
  equals the law of $(X_t)_t$. Then every shift preserves the law, the marginals are identically
  distributed, sliding-window functions of a stationary sequence are stationary, and iid sequences
  are stationary.
\end{definition}

\begin{definition}[Conditions $D(u_n)$ and $D'(u_n)$]
  \label{def:leadbetter}
  \uses{def:stationary}
  \lean{ExtremeValueTheory.partialMax, ExtremeValueTheory.maxLE, ExtremeValueTheory.ConditionD, ExtremeValueTheory.ConditionD'}
  \leanok
  Let $M_n = \max_{t < n} X_t$. Condition $D(u_n)$ holds if there are $l_n = o(n)$ and
  $\alpha_n \to 0$ such that for all $A, B \subseteq \{0, \dots, n-1\}$ separated by at least $l_n$,
  \begin{align*}
    \big| \PR[\max_{A \cup B} X \le u_n] - \PR[\max_A X \le u_n] \, \PR[\max_B X \le u_n] \big|
      \; \le \; \alpha_n .
  \end{align*}
  Condition $D'(u_n)$ holds if
  $\limsup_n n \sum_{1 \le j < n/k} \PR[X_0 > u_n, X_j > u_n] \to 0$ as $k \to \infty$
  [LLR, \S3.2, (3.4.3)].
\end{definition}

\begin{theorem}[Block approximation]
  \label{thm:block-approx}
  \uses{def:leadbetter}
  \lean{ProbabilityTheory.IsStationary.abs_measureReal_maxLE_range_sub_pow_le,
    ExtremeValueTheory.tendsto_measureReal_partialMax_sub_pow, MeasureTheory.sum_measureReal_sub_sum_inter_le_measureReal_biUnion}
  \leanok
  For a stationary sequence satisfying $D(u_n)$ with $n \PR[X_0 > u_n] \to \tau$, and fixed
  $k \ge 1$, $\PR[M_n \le u_n] - \PR[M_{\lfloor n/k \rfloor} \le u_n]^k \to 0$ [LLR, Lemma~3.3.1].
\end{theorem}
\begin{proof}
  \leanok
  Split $\{0, \dots, n-1\}$ into $k$ blocks separated by gaps of length $l_n$. The gaps cost at most
  $2 k l_n \PR[X_0 > u_n] \to 0$, and $D(u_n)$ factorizes the blocks with error $k \alpha_n$.
\end{proof}

\begin{theorem}[Leadbetter's theorem]
  \label{thm:leadbetter}
  \uses{thm:block-approx}
  \lean{ExtremeValueTheory.tendsto_measureReal_partialMax_le_exp, ProbabilityTheory.IsStationary.measureReal_maxLE_range_bounds}
  \leanok
  If a stationary sequence satisfies $D(u_n)$ and $D'(u_n)$ and $n \PR[X_0 > u_n] \to \tau$, then
  $\PR[M_n \le u_n] \to e^{-\tau}$, as for iid sequences [LLR, Theorem~3.4.1], [EKM, Theorem~4.4.3].
\end{theorem}
\begin{proof}
  \leanok
  By Bonferroni, $1 - \frac{n}{k} p_n \le \PR[M_{n/k} \le u_n] \le 1 - \frac{n}{k} p_n + \frac{n}{k}
  \sum_{1 \le j < n/k} \PR[X_0 > u_n, X_j > u_n]$ with $p_n = \PR[X_0 > u_n]$. Raise to the $k$-th
  power, use Theorem~\ref{thm:block-approx}, and let $k \to \infty$ using $D'(u_n)$.
\end{proof}

\section{The extremal index}

\begin{definition}[Extremal index]
  \label{def:extremal-index}
  \uses{def:leadbetter}
  \lean{ExtremeValueTheory.HasExtremalIndex}
  \leanok
  A stationary sequence has extremal index $\theta \in [0, 1]$ if for every $\tau > 0$ and every
  sequence $u_n$ with $n \PR[X_0 > u_n] \to \tau$, $\PR[M_n \le u_n] \to e^{-\theta\tau}$
  [LLR, \S3.7], [EKM, Definition~8.1.1].
\end{definition}

\begin{theorem}[Maxima of stationary sequences]
  \label{thm:stationary-max}
  \uses{def:extremal-index, thm:leadbetter, thm:three-types-of-extr-val-distr}
  \lean{ExtremeValueTheory.hasExtremalIndex_one, ExtremeValueTheory.tendsto_measureReal_partialMax_le_pow,
    ExtremeValueTheory.isExtremeValueDistr_pow, CumulativeDistributionFunction.IsExtremeValueDistr.continuous}
  \leanok
  Under $D$ and $D'$ at all levels, the extremal index is $1$. If the marginal $F$ satisfies
  $F^n(a_n x + b_n) \to H(x)$ for an extreme value distribution $H$ and the sequence has extremal
  index $\theta > 0$, then $\PR[M_n \le a_n x + b_n] \to H(x)^\theta$, and $H^\theta$ is an extreme
  value distribution of the same type as $H$ [LLR, Theorem~3.7.1], [EKM, Theorem~8.1.2].
\end{theorem}
\begin{proof}
  \leanok
  $n(1 - F(a_n x + b_n)) \to -\log H(x)$, so the definition of $\theta$ applies with
  $\tau = -\log H(x)$. For the type, $H^\theta = H^{\theta}$ is max-stable, and the
  classification gives an affine map with $A \cdot H = H^\theta$.
\end{proof}

\begin{proposition}[Examples]
  \label{prop:extremal-index-examples}
  \uses{thm:stationary-max}
  \lean{ExtremeValueTheory.measureReal_maxLE_of_iIndepFun, ExtremeValueTheory.conditionD_of_iIndepFun,
    ExtremeValueTheory.conditionD'_of_iIndepFun, ExtremeValueTheory.hasExtremalIndex_one_of_iIndepFun, ExtremeValueTheory.isStationary_movingMax,
    ExtremeValueTheory.hasExtremalIndex_movingMax}
  \leanok
  Iid sequences satisfy $D$ and $D'$ and have extremal index $1$. The moving maximum
  $X_t = \max(Z_t, Z_{t+1})$ of an iid sequence is stationary with extremal index $1/2$
  [EKM, Example~8.1.3].
\end{proposition}
\begin{proof}
  \leanok
  For the moving maximum, $\PR[M_n \le u] = G(u)^{n+1}$ with $G$ the law of $Z_0$, while
  $\PR[X_0 > u] \sim 2(1 - G(u))$.
\end{proof}

\section{Regularly varying time series and the tail process}

We follow Basrak and Segers, \emph{Regularly varying multivariate time series}, Stoch.\ Proc.\ Appl.119 (2009) [BS09], Planini\'c and Soulier, \emph{The tail process revisited}, Extremes 21 (2018)
[PS18], and Kulik and Soulier, \emph{Heavy-Tailed Time Series}, Springer (2020) [KS20].

\begin{definition}[Conditional limits, tail process and spectral tail process]
  \label{def:tail-process}
  \uses{def:regularly-varying}
  \lean{ExtremeValueTheory.TendstoLawCond, ExtremeValueTheory.HasSpectralTailProcess, ExtremeValueTheory.HasTailProcess,
    ExtremeValueTheory.HasShiftInvariantWindows}
  \leanok
  A sequence $(X_t)_{t \in \bZ}$ has spectral tail process $(\Theta_t)$ with index $\alpha > 0$ if
  $\PR[|X_0| > x]$ is regularly varying with index $-\alpha$ and, for every finite window
  $S \subset \bZ$, the conditional law of $(X_t/|X_0|)_{t \in S}$ given $|X_0| > x$ converges weakly to
  the law of $(\Theta_t)_{t \in S}$ as $x \to \infty$. It has tail process $(Y_t)$ if $|Y_0|$ is
  Pareto($\alpha$) and the conditional law of $(X_t/x)_{t\in S}$ given $|X_0| > x$ converges to that of
  $(Y_t)_{t \in S}$ [BS09, \S2], [KS20, Ch.~5].
\end{definition}

\begin{theorem}[Polar decomposition (Basrak--Segers)]
  \label{thm:basrak-segers}
  \uses{def:tail-process}
  \lean{ExtremeValueTheory.tendstoLawCond_polar, ExtremeValueTheory.isRegularlyVarying_of_tendstoLawCond,
    ExtremeValueTheory.tendstoLawCond_spectral_of_tendstoLawCond, ExtremeValueTheory.eq_map_prod_of_tendstoLawCond, ExtremeValueTheory.polar_iff,
    ExtremeValueTheory.HasSpectralTailProcess.hasTailProcess, ExtremeValueTheory.HasTailProcess.hasSpectralTailProcess,
    ExtremeValueTheory.HasTailProcess.map_window_eq, ExtremeValueTheory.HasSpectralTailProcess.ae_abs_eq_one}
  \leanok
  A sequence has a spectral tail process if and only if it has a tail process, and then
  $Y_t = |Y_0| \Theta_t$ with $|Y_0|$ Pareto($\alpha$) independent of $(\Theta_t)$, and $|\Theta_0| = 1$
  [BS09, Theorem~2.1].
\end{theorem}
\begin{proof}
  \leanok
  Regular variation makes $|X_0|/x$ asymptotically Pareto given $|X_0| > x$; conditioning further on
  the radial part shows that it decouples from the angular part $X/|X_0|$ in the limit.
\end{proof}

\begin{theorem}[Time-change formula]
  \label{thm:time-change}
  \uses{thm:basrak-segers, def:stationary}
  \lean{ExtremeValueTheory.HasSpectralTailProcess.lintegral_timeChange}
  \leanok
  For a stationary sequence with spectral tail process $\Theta$, all $i \in \bZ$, finite windows
  $S$ and measurable $g \ge 0$,
  \begin{align*}
    \E\big[ g((\Theta_{t-i})_{t\in S}) \, \mathbf 1\{\Theta_{-i} \ne 0\} \big]
      \; = \; \E\big[ g((\Theta_t / |\Theta_i|)_{t \in S}) \, |\Theta_i|^\alpha \big]
  \end{align*}
  [BS09, Theorem~3.1], [KS20, Theorem~5.3.1].
\end{theorem}
\begin{proof}
  \leanok
  Stationarity moves the conditioning from $|X_0| > x$ to $|X_i| > x$. Comparing the two
  conditionings in the tail-process limit at every level $\varepsilon > 0$ gives a measure identity,
  which is then rewritten in terms of $\Theta$.
\end{proof}

\begin{definition}[Candidate extremal index]
  \label{def:candidate-extremal-index}
  \uses{def:tail-process}
  \lean{ExtremeValueTheory.candidateExtremalIndex}
  \leanok
  $\vartheta = \E\big[\sup_{k \ge 0} |\Theta_k|^\alpha - \sup_{k \ge 1} |\Theta_k|^\alpha\big]$, with
  suprema in $[0, \infty]$ and $\infty - \infty = 0$ [KS20, \S7.5], [PS18].
\end{definition}

\begin{theorem}[Properties of the candidate extremal index]
  \label{thm:candidate-extremal-index}
  \uses{def:candidate-extremal-index, thm:time-change}
  \lean{ExtremeValueTheory.candidateExtremalIndex_le_one, ExtremeValueTheory.candidateExtremalIndex_eq_measure,
    ExtremeValueTheory.HasTailProcess.candidateExtremalIndex_eq, ExtremeValueTheory.HasSpectralTailProcess.ae_eq_zero_of_indepFun,
    ExtremeValueTheory.HasSpectralTailProcess.candidateExtremalIndex_eq_one_of_indepFun,
    ExtremeValueTheory.HasSpectralTailProcess.candidateExtremalIndex_pos}
  \leanok
  $\vartheta \le 1$ and $\vartheta = \PR[\sup_{k \ge 1} |Y_k| \le 1]$. If $X_t$ is independent of $X_0$
  then $\Theta_t = 0$; in particular $\vartheta = 1$ for iid sequences. For a stationary sequence
  with $\Theta_t \to 0$ as $t \to \infty$, $\vartheta > 0$ [PS18, \S4], [KS20, \S7.5].
\end{theorem}
\begin{proof}
  \leanok
  $|\Theta_0| = 1$ gives $\vartheta \le 1$. Writing $Y = R\Theta$ with $R$ Pareto independent of
  $\Theta$, $\PR[\sup_{k\ge1}|Y_k| \le 1 \mid \Theta] = \sup_{k\ge0}|\Theta_k|^\alpha -
  \sup_{k\ge1}|\Theta_k|^\alpha$. Positivity: decompose according to the last time $|Y_t| > 1$, which
  exists since $\Theta_t \to 0$, and apply the time-change formula to each piece.
\end{proof}

\section{The extremal index and the tail process}

\begin{lemma}[Stationarity notions agree]
  \label{lem:stationary-windows}
  \uses{def:stationary, def:tail-process}
  \lean{ExtremeValueTheory.hasShiftInvariantWindows_iff_isStationary, ExtremeValueTheory.HasShiftInvariantWindows.isStationary_nat,
    ExtremeValueTheory.HasShiftInvariantWindows.map_comp}
  \leanok
  For measurable $(X_t)_{t \in \bZ}$ under a finite measure, shift-invariance of all finite windows is
  equivalent to stationarity of every shifted $\bN$-indexed sequence $(X_{n+k})_{n \in \bN}$.
\end{lemma}
\begin{proof}
  \leanok
  Finite-dimensional laws determine the law of the path.
\end{proof}

\begin{definition}[Anticlustering and condition $\mathcal A$]
  \label{def:anticlustering}
  \uses{def:tail-process}
  \lean{ExtremeValueTheory.AntiClusteringAt, ExtremeValueTheory.ForwardAntiClusteringAt, ExtremeValueTheory.AntiClustering, ExtremeValueTheory.ConditionA}
  \leanok
  $AC(r_n)$ at levels $u_n$: $\lim_m \limsup_n \PR[\max_{m \le |t| \le r_n} |X_t| > u_n \mid |X_0| > u_n] = 0$
  (the forward version only uses $m \le t \le r_n$). Condition $\mathcal A(u_n)$:
  $\PR[M_n \le u_n] - \PR[M_{r_n} \le u_n]^{\lfloor n/r_n \rfloor} \to 0$ [BS09, \S4], [KS20, \S6.1].
\end{definition}

\begin{theorem}[Anticlustering forces $\Theta_t \to 0$]
  \label{thm:ac-theta-zero}
  \uses{def:anticlustering, thm:candidate-extremal-index}
  \lean{ExtremeValueTheory.HasSpectralTailProcess.tendsto_zero_of_forwardAntiClusteringAt,
    ExtremeValueTheory.HasSpectralTailProcess.candidateExtremalIndex_pos_of_forwardAntiClusteringAt,
    ExtremeValueTheory.HasSpectralTailProcess.candidateExtremalIndex_pos_of_antiClustering}
  \leanok
  Forward anticlustering at a single level sequence $u_n \to \infty$ with $r_n \to \infty$ implies
  $\Theta_t \to 0$ almost surely as $t \to \infty$; neither stationarity nor $r_n/n \to 0$ is used.
  For stationary sequences it follows that $\vartheta > 0$ [BS09, Proposition~4.2].
\end{theorem}
\begin{proof}
  \leanok
  The tail-process limit of the anticlustering probabilities at level $1$, extended to all levels
  through the independent Pareto radius $Y = R\Theta$.
\end{proof}

\begin{theorem}[The extremal index equals the candidate extremal index]
  \label{thm:theta-eq-vartheta}
  \uses{thm:ac-theta-zero, thm:time-change, def:extremal-index, lem:stationary-windows}
  \lean{ExtremeValueTheory.HasShiftInvariantWindows.measureReal_exists_lt_eq_sum, ExtremeValueTheory.HasSpectralTailProcess.tendsto_blockRatio,
    ExtremeValueTheory.HasSpectralTailProcess.tendsto_partialMax_abs, ExtremeValueTheory.HasSpectralTailProcess.hasExtremalIndex_abs, ExtremeValueTheory.HasSpectralTailProcess.hasExtremalIndex_of_nonneg,
    ExtremeValueTheory.HasSpectralTailProcess.tendsto_partialMax_abs_of_antiClustering}
  \leanok
  For a stationary sequence with spectral tail process $\Theta$, $u_n \to \infty$, $r_n \to \infty$
  and forward anticlustering at $u_n$,
  \begin{align*}
    \frac{\PR[\max_{t < r_n} |X_t| > u_n]}{r_n \PR[|X_0| > u_n]} \; \to \; \vartheta ;
  \end{align*}
  no mixing and no condition $r_n/n \to 0$ is needed for this. If in addition
  $n \PR[|X_0| > u_n] \to \tau$, $r_n/n \to 0$ and $\mathcal A(u_n)$ holds, then
  $\PR[\max_{t<n} |X_t| \le u_n] \to e^{-\vartheta\tau}$. Under $AC(r_n)$ and $\mathcal A(a_n x)$ at all
  levels, $|X|$ (and $X$ itself if $X \ge 0$) has extremal index $\vartheta$, and
  $\PR[M_n \le a_n x] \to \exp(-\vartheta x^{-\alpha})$ [BS09, Proposition~4.2], [KS20, \S7.5].
\end{theorem}
\begin{proof}
  \leanok
  Exact first-exceedance decomposition
  $\PR[M_r > x] = \sum_{k<r} \PR[|X_0| > x, \max_{0 < t \le k} |X_t| \le x]$ by stationarity; each term,
  divided by $\PR[|X_0| > x]$, tends to $\PR[\max_{0<t\le k} |Y_t| \le 1]$, and anticlustering controls
  the terms with $k \ge m$. The limit of the Ces`aro average is $\vartheta$. Condition $\mathcal A$
  then gives $(1 - r_n \vartheta \tau / n)^{n/r_n} \to e^{-\vartheta\tau}$.
\end{proof}

\begin{theorem}[$\theta \le \vartheta$ without anticlustering]
  \label{thm:theta-le-vartheta}
  \uses{thm:theta-eq-vartheta}
  \lean{ExtremeValueTheory.HasSpectralTailProcess.le_candidateExtremalIndex_of_tendsto_blockRatio, ExtremeValueTheory.HasSpectralTailProcess.le_candidateExtremalIndex_of_tendsto,
    ExtremeValueTheory.HasSpectralTailProcess.le_candidateExtremalIndex, ExtremeValueTheory.tendsto_measureReal_forall_le_candidateExtremalIndex}
  \leanok
  Every limit point of the cluster ratio is at most $\vartheta$, using only stationarity, the tail
  process, $u_n \to \infty$ and $r_n \to \infty$. Consequently, if $|X|$ has extremal index $\theta$ and
  condition $\mathcal A(u_n)$ holds along one level sequence with $n \PR[|X_0| > u_n] \to \tau > 0$,
  $r_n \to \infty$ and $r_n/n \to 0$, then $\theta \le \vartheta$. Anticlustering is not needed. (The
  hypotheses of the literature version, [KS20, Lemma~7.5.4], were not checked against the source.)
\end{theorem}
\begin{proof}
  \leanok
  Without anticlustering, the first-exceedance terms with $k \ge m$ are bounded above by their limits
  but may lose mass, which gives the inequality.
\end{proof}

\begin{proposition}[Moving maximum: $\theta = \vartheta = 1/2$]
  \label{prop:moving-max-vartheta}
  \uses{thm:candidate-extremal-index, prop:extremal-index-examples}
  \lean{ExtremeValueTheory.movingMax, ExtremeValueTheory.isRegularlyVarying_movingMax, ExtremeValueTheory.hasSpectralTailProcess_movingMax,
    ExtremeValueTheory.candidateExtremalIndex_movingMax, ExtremeValueTheory.hasExtremalIndex_movingMax_candidateExtremalIndex}
  \leanok
  For iid Pareto($\alpha$) $Z$ and $X_t = \max(Z_t, Z_{t+1})$, $\PR[X_0 > x] = 2x^{-\alpha} - x^{-2\alpha}$,
  the spectral tail process is the indicator of $\{-1, 0\}$ or of $\{0, 1\}$ with probability $1/2$
  each, and $\vartheta = 1/2$, matching the directly computed extremal index $\theta = 1/2$.
\end{proposition}
\begin{proof}
  \leanok
  An exceedance of $X_0$ comes from $Z_0$ or $Z_1$, each with conditional probability $1/2$ in the limit.
\end{proof}
